% original_pdf: source/普通高考/2012/2012福建文.pdf
% page: 1 (question 6)
% embedded_bitmap_attempt: pdfimages -list returned no embedded images; redraw from rendered page
% evidence: tmp/review_flowchart_2012_fujian_liberal/grid/q6_grid100.jpg;
%   q6_original_final.png (no-grid crop from the source page)
% elements: rounded 开始/结束 terminals; initialization rectangle k=1,s=1; condition diamond k<4?;
% process rectangles s=2s−k and k=k+1; output parallelogram 输出 s; solid arrows and 是/否 labels.
% node-edge list: 开始→(k=1,s=1)→k<4?; 是→s=2s−k→k=k+1→left loop-back to the connector above
% k<4?; 否→输出 s→结束. The output branch is independent and the main down arrow is unique.
% solid segments: all node borders and directed connectors; dashed segments: none.
% Branch labels use natural edge-relative anchors; node text is locally zero-indented and centered.

\begin{tikzpicture}[
  baseline,
  >=stealth,
  x=1cm,
  y=1cm,
  line width=0.55pt,
  line cap=round,
  line join=round,
  font=\normalsize,
  flowtext/.style={anchor=center,align=center,
    font=\normalsize\setlength{\parindent}{0pt}\noindent},
  every node/.style={flowtext},
  flowbox/.style={flowtext,draw,minimum height=0.68cm,inner xsep=4pt,inner ysep=2.5pt},
  branchlabel/.style={flowtext,inner xsep=1pt,inner ysep=1pt},
  terminal/.style={flowbox,rounded corners=0.20cm,minimum width=1.24cm},
  io/.style={flowbox,shape=trapezium,trapezium left angle=72,trapezium right angle=108,
    minimum height=0.76cm},
  decision/.style={flowbox,shape=diamond,aspect=2.8,inner xsep=4pt,inner ysep=2pt,
    minimum height=1.00cm}
]
  % Main vertical path.
  \node[terminal] (start) at (0,10.00) {开始};
  \node[flowbox,minimum width=2.85cm] (init) at (0,8.55) {$k=1,s=1$};
  \node[decision,minimum width=2.90cm] (test) at (0,4.10) {$k<4?$};
  \node[flowbox,minimum width=2.65cm] (update) at (2.60,5.30) {$s=2s-k$};
  \node[flowbox,minimum width=2.35cm] (increment) at (2.60,6.80) {$k=k+1$};
  \node[io,minimum width=2.15cm] (output) at (0,2.10) {输出 $s$};
  \node[terminal] (finish) at (0,0.60) {结束};

  \draw[->] (start.south) -- (init.north);
  % Split the stem at the loop-entry level so the return arrow has its own
  % left-pointing head and the main path has one downward arrow into test.
  \coordinate (merge-level) at (0,7.65);
  \draw (init.south) -- (merge-level);
  \draw[->] (merge-level) -- (test.north);
  \draw[->] (test.south) -- node[branchlabel,right] {否} (output.north);
  \draw[->] (output.south) -- (finish.north);

  % Yes branch rises through the two right-hand updates and returns left.
  \draw[->] (test.east) -- node[branchlabel,above] {是} (2.60,4.10) -- (update.south);
  \draw[->] (update.north) -- (increment.south);
  \draw[->] (increment.north) -- (2.60,7.65) -- (merge-level);
  \path[use as bounding box] (-1.70,10.35) rectangle (4.00,0.22);
\end{tikzpicture}
